% !TeX root = main.tex
\chapter{Implémentation}
\label{ch:implementation}

Ce chapitre décrit la couche \texttt{streamkm.spark}, opérateur par
opérateur, en suivant le même plan que le chapitre~4 de la thèse : d'abord
le dataflow borné (§\ref{sec:impl-dataflow}, miroir du §4.1 de la thèse),
puis la méthodologie d'évaluation du coût (§\ref{sec:impl-evaluation},
miroir du §4.2), puis la variante à requêtes périodiques
(§\ref{sec:impl-requetes}, miroir du §4.3, \emph{décrite} plutôt
qu'implémentée --- cf. justification en fin de section), et enfin un bilan
de ce qui se traduit ou non depuis Flink (§\ref{sec:portage-flink-spark}).
Rappel de la règle d'or (chapitre~\ref{ch:introduction}) : aucun des
fichiers décrits ici ne contient de logique algorithmique propre --- tout le
calcul vient de \texttt{streamkm.core} et \texttt{streamkm.coreset}
(chapitre~\ref{ch:algorithmes}), déjà testés indépendamment de Spark.

\section{StreamKM++ distribué, flux borné}
\label{sec:impl-dataflow}

\subsection{Coresets partiels par partition}

\texttt{StreamKMPlusPlus.partialCoresets} porte l'opérateur
\texttt{FlatMapToPartialCoresets} de la thèse (§4.1) : chaque subtask Flink
« consomme les points d'entrée selon l'algorithme 2.5 » --- une partition
Spark joue exactement ce rôle.

\begin{verbatim}
def partialCoresets(rdd: RDD[Point], p: StreamKMParams): RDD[Array[Point]] =
  rdd.mapPartitionsWithIndex { (idx, it) =>
    val bm = new BucketManager(p.m, p.seed + idx)
    bm.insertAll(it)
    Iterator.single(bm.extractCoreset())
  }
\end{verbatim}

La graine de chaque \texttt{BucketManager} est dérivée de l'indice de
partition (\texttt{p.seed + idx}) plutôt que partagée entre partitions ---
condition nécessaire pour que l'expérience E3bis
(chapitre~\ref{ch:evaluation}) reste interprétable : avec une graine
partagée, changer le nombre de partitions changerait aussi la suite
aléatoire consommée par chaque partition, confondant l'effet mesuré (le
partitionnement) avec un simple changement de graine.

\subsection{Fusion arborescente du coreset final}

\texttt{StreamKMPlusPlus.mergeCoresets} remplace l'opérateur
\texttt{FlatMapToFinalCoreset} de la thèse --- l'opérateur \textbf{non
parallèle} par construction (§4.1 : fusion séquentielle, un coreset partiel
à la fois) que la thèse elle-même identifie comme le goulot qui plafonne le
speedup au-delà de 8 c\oe{}urs (§5.3.1). C'est notre apport principal par
rapport à la thèse, rendu possible par le découplage partitions/c\oe{}urs de
Spark (§\ref{sec:spark-treereduce}) :

\begin{verbatim}
def mergeCoresets(rdd: RDD[Point], p: StreamKMParams, depth: Int = 2): Array[Point] = {
  val perPartition: RDD[BucketManager] = rdd.mapPartitionsWithIndex { (idx, it) =>
    val bm = new BucketManager(p.m, p.seed + idx)
    bm.insertAll(it)
    Iterator.single(bm)
  }
  perPartition.treeReduce((a, b) => a.merge(b), depth).extractCoreset()
}
\end{verbatim}

Le \texttt{combOp} de ce \texttt{treeReduce} est \texttt{BucketManager.merge}
(chapitre~\ref{ch:algorithmes}, §\ref{sec:algo-merge-reduce}), réutilisé
\emph{tel quel} --- aucune extension de \texttt{streamkm.coreset} n'a été
nécessaire pour distribuer merge-and-reduce sur un cluster, ce qui confirme
a posteriori la séparation stricte des couches annoncée en introduction.

\subsection{Boucle Structured Streaming complète}

\texttt{StreamKMPlusPlus.run} assemble l'ensemble du dataflow --- de
l'équivalent de \texttt{HdfsSource} jusqu'à l'équivalent de
\texttt{FlatMapToKmeansPP} --- autour d'un état global unique
(\texttt{BucketManager}) maintenu côté \emph{driver} entre les micro-batches :

\begin{verbatim}
def run(source: DataFrame, p: StreamKMParams,
        onModel: KMeansModel => Unit,
        featuresCol: String = "features"): StreamingQuery = {

  val global = new BucketManager(p.m, p.seed)

  source.writeStream.foreachBatch { (batch, batchId) =>
    val points = batch.rdd.map(row => Point(row.getAs[Seq[Double]](featuresCol).toArray))
    val partial = mergeCoresets(points, p.copy(seed = p.seed + batchId))
    partial.foreach(global.insert)
    val model = KMeans.fit(global.extractCoreset(), p.k, p.nRestarts, seed = p.seed)
    onModel(model)
  }.start()
}
\end{verbatim}

Chaque micro-batch produit son propre coreset distribué (\texttt{mergeCoresets}),
qui est ensuite intégré dans l'état global par \texttt{insert} --- exactement
la même opération de fusion que \texttt{FlatMapToFinalCoreset}, appliquée
cette fois entre micro-batches plutôt qu'entre partitions. Le contrat
d'entrée (une colonne \texttt{featuresCol} de type \texttt{array<double>},
poids implicite $1.0$) reflète directement \texttt{Point.scala}
(chapitre~\ref{ch:algorithmes}) : le poids n'existe que dans les coresets
produits en interne, jamais dans les données brutes du flux.

\subsection{Table de correspondance}

\begin{table}[h]
\centering
\small
\begin{tabular}{p{3.6cm}p{4.6cm}p{4.6cm}}
\toprule
\textbf{Opérateur Flink (thèse)} & \textbf{Rôle} & \textbf{Équivalent Spark} \\
\midrule
\texttt{HdfsSource} & lecture parallèle des lignes & source fichier partitionnée \\
\texttt{HandleWatermark} & détecter la fin de fichier & \textbf{aucun} --- \texttt{foreachBatch} rend la frontière explicite \\
\texttt{FlatMapToPoint} & parsing \texttt{String} $\to$ \texttt{Point} & \texttt{mapPartitions} + \texttt{filter} \\
\texttt{FlatMapToPartialCoresets} & merge-and-reduce par subtask & \texttt{partialCoresets} \\
\texttt{FlatMapToFinalCoreset} & fusion séquentielle (par. 1) & \texttt{mergeCoresets} (\texttt{treeReduce}) \\
\texttt{FlatMapToKmeansPP} & 5$\times$ $k$-means++ (par. 1) & \texttt{KMeans.fit} sur le driver \\
\texttt{HdfsSink} & écriture des centres & sink (fichier ou \texttt{onModel}) \\
\bottomrule
\end{tabular}
\caption{Correspondance des opérateurs, dataflow borné (§4.1 de la thèse).}
\end{table}

\section{Méthodologie d'évaluation du coût}
\label{sec:impl-evaluation}

La thèse (§4.2) recalcule le coût d'un clustering en parallèle, à partir
d'un fichier de centres déjà connu, dans un job Flink séparé --- nécessaire
pour mesurer la qualité sur de gros volumes sans que ce calcul devienne
lui-même le goulot d'étranglement des expériences du chapitre~5. Elle utilise
pour cela un \emph{connected stream} (deux sources fusionnées : les points
et les centres), un opérateur de diffusion et un opérateur de somme partielle.
Cette machinerie tient, côté Spark, en trois lignes :

\begin{verbatim}
def cost(points: RDD[Point], centers: Array[Array[Double]], depth: Int = 2): Double = {
  val bcCenters = points.sparkContext.broadcast(centers)
  points
    .mapPartitions(it => Iterator.single(KMeans.cost(it.toArray, bcCenters.value)))
    .treeReduce(_ + _, depth)
}
\end{verbatim}

\texttt{sc.broadcast(centers)} remplace \texttt{FlatMapToCentroid} +
\texttt{.broadcast()} ; \texttt{mapPartitions(... KMeans.cost ...)} remplace
\texttt{CoFlatMapToPartialCost} ; \texttt{treeReduce(\_+\_)} remplace
\texttt{CoFlatMapToFinalCost}. Contrairement à la fusion de coresets
(§\ref{sec:impl-dataflow}), cette agrégation est une simple somme ---
associative et exacte à la précision flottante près --- ce qui explique
pourquoi elle ne nécessite aucune des précautions de reproductibilité par
partition discutées plus haut : le résultat ne dépend pas du nombre de
partitions choisi. C'est cet instrument qui sert de mesure à
\textbf{toutes} les expériences du chapitre~\ref{ch:evaluation} --- jamais
le coût mesuré sur le coreset lui-même, qui n'est pas comparable au coût sur
les données complètes.

\section{Variante à requêtes périodiques (décrite, non implémentée)}
\label{sec:impl-requetes}

Le §4.3 de la thèse propose une variante où le clustering est ré-évalué
périodiquement plutôt qu'une seule fois à la fin du flux, et où plusieurs
valeurs de $k$ sont calculées simultanément. Faute de temps avant la
deadline, cette variante \textbf{n'a pas été implémentée} --- conformément à
notre méthodologie de travail, qui préférait explicitement documenter cette
limite plutôt que de livrer une implémentation bâclée. La thèse elle-même ne
l'évalue pas expérimentalement (elle le signale comme travail futur, ch.~6),
ce qui rend ce choix d'autant plus défendable.

Le portage resterait néanmoins direct avec ce que nous avons déjà construit :
le \emph{trigger} de Structured Streaming remplacerait directement
l'opérateur \texttt{SourceRequests\_R} de génération de requêtes périodiques
(\texttt{Trigger.ProcessingTime("30 seconds")} au lieu d'un flux de requêtes
explicite) ; \texttt{BucketManager.extractCoreset} --- qui ne vide jamais
l'état, une propriété déjà exploitée par \texttt{run} --- permettrait de
dupliquer le coreset courant pour plusieurs valeurs de $k$ simultanément
(\texttt{Seq(25, 50, 75).map(k => KMeans.fit(coreset, k))}), remplaçant
\texttt{FlatMapToFinalCoreset\_R}. Seule la persistance du meilleur
clustering \emph{par} $k$ (\texttt{ValueState} chez Flink) demanderait une
structure nouvelle côté Spark (une simple \texttt{Map[Int, KMeansModel]}
tenue par le \texttt{onModel} de \texttt{run} suffirait).

Le §4.3.1 de la thèse (\emph{Queryable State}, interrogation en temps réel de
l'état interne depuis l'extérieur du job) \textbf{n'a en revanche pas
d'équivalent Spark} : Spark n'expose aucun état inter-micro-batch
interrogeable depuis l'extérieur. C'est un avantage architectural réel de
Flink, que nous assumons comme tel --- voir la discussion du
chapitre~\ref{ch:discussion}. Le contournement usuel (un sink \texttt{memory}
interrogé par \texttt{spark.sql("SELECT * FROM centers")}, ou l'écriture des
centres dans un store externe) n'offre pas la même garantie de fraîcheur
immédiate.

\section{Ce qui se traduit depuis Flink, ce qui ne se traduit pas}
\label{sec:portage-flink-spark}

Ce paragraphe de synthèse répond directement à la question posée par
l'énoncé --- porter en Spark une solution Flink --- et n'a pas
d'équivalent dans la thèse : c'est notre contribution propre.

\begin{itemize}
  \item \textbf{Disparaît entièrement} : \texttt{HandleWatermark}. Le
    bricolage watermark/marqueur « EOF » de la thèse n'a plus de raison
    d'être --- \texttt{foreachBatch} rend la frontière de micro-batch
    explicite par construction (§\ref{ch:spark}).
  \item \textbf{Se traduit avec un gain net} : la fusion des coresets
    partiels. L'opérateur non-parallèle de la thèse (\S\ref{sec:impl-dataflow})
    devient une réduction en arbre (\texttt{treeReduce}), ce qui repousse
    directement le plafond de speedup identifié au §5.3.1 de la thèse ---
    à confirmer expérimentalement (chapitre~\ref{ch:evaluation},
    expérience E7).
  \item \textbf{Se traduit à l'identique, en plus concis} : la méthodologie
    d'évaluation du coût (§\ref{sec:impl-evaluation}) --- trois opérateurs
    Flink dédiés et un \emph{connected stream} deviennent trois lignes de
    \texttt{broadcast} + \texttt{treeReduce}.
  \item \textbf{Se traduit, avec un degré de liberté nouveau} : le
    parallélisme. Le couplage partitions/c\oe{}urs de Flink se découple côté
    Spark (§\ref{sec:spark-partitions-coeurs}), ce qui permet une expérience
    (E3bis) que le modèle de Flink ne permet pas.
  \item \textbf{Ne se traduit pas} : Queryable State
    (§\ref{sec:impl-requetes}). C'est la seule perte de fonctionnalité nette
    identifiée dans ce portage --- discutée au chapitre~\ref{ch:discussion}.
\end{itemize}
